\begin{figure}[H]
\centering
\begin{tikzpicture}[
    x={(2.1cm,0cm)}, y={(1.0cm,0.8cm)}, z={(0cm,2.0cm)},
    >=Stealth, font=\small]

  % plano degenerado E_1
  \fill[azulUNI!10] (-1.1,-1.1,0)--(1.1,-1.1,0)--(1.1,1.1,0)--(-1.1,1.1,0)--cycle;
  \draw[azulUNI!45] (-1.1,-1.1,0)--(1.1,-1.1,0)--(1.1,1.1,0)--(-1.1,1.1,0)--cycle;
  \node[azulUNI,anchor=south east,font=\small] at (-0.85,1.1,0)
        {$E_1$ \; ($A=1$, degenerado)};

  % base canónica (no es propia de B)
  \draw[->,thick,grisUNI] (0,0,0)--(1,0,0) node[below right]{$|e_1\rangle$};
  \draw[->,thick,grisUNI] (0,0,0)--(0,1,0) node[above left]{$|e_2\rangle$};

  % base propia común
  \draw[->,very thick,rojoUNI]  (0,0,0)--(0.85,0.85,0)
        node[above right]{$|u_1\rangle$};
  \draw[->,very thick,azulUNI]  (0,0,0)--(0.85,-0.85,0)
        node[below]{$|u_2\rangle$};
  \draw[->,very thick,verdeUNI] (0,0,0)--(0,0,1.4)
        node[above]{$|v_3\rangle$};

  % B mezcla e_1 y e_2
  \draw[<->,dashed,grisUNI] (0.55,0,0) to[bend left=35] (0,0.55,0);

  % leyenda
  \node[align=left,anchor=north west,font=\small] at (canvas cs:x=-3.6cm,y=-1.7cm) {%
    \textcolor{rojoUNI}{$|u_1\rangle=\tfrac{1}{\sqrt2}(|e_1\rangle+|e_2\rangle)$: $A=1,\;B=+2$}\\[2pt]
    \textcolor{azulUNI}{$|u_2\rangle=\tfrac{1}{\sqrt2}(|e_1\rangle-|e_2\rangle)$: $A=1,\;B=-2$}\\[2pt]
    \textcolor{verdeUNI}{$|v_3\rangle=|e_3\rangle$: $A=-2,\;B=3$}\\[6pt]
    \textcolor{grisUNI}{$B|e_1\rangle=2|e_2\rangle$\; (\,$B$ «mezcla» dentro de $E_1$)}};
\end{tikzpicture}
\caption{Diagonalización simultánea con degeneración (Ejemplo de la Sección~2.8.3). Dentro del
plano $E_1$, $A$ es un múltiplo de la identidad, pero $B$ gira los vectores:
la base canónica $\{|e_1\rangle,|e_2\rangle\}$ no es propia de $B$. La base
común se obtiene diagonalizando $B$ restringido a $E_1$, lo que equivale a
girar $45^\circ$ dentro del plano; $|v_3\rangle$ ya era propio de ambos.}
\label{fig:diag-simultanea}
\end{figure}
